\documentclass[11pt,a4paper]{article}
\usepackage[a4paper,margin=24mm]{geometry}
\usepackage{amsmath,amssymb,bm}
\usepackage{booktabs,longtable,array}
\usepackage{xcolor,listings}
\usepackage[hidelinks]{hyperref}
\definecolor{codebg}{HTML}{F4F6F9}
\definecolor{codeframe}{HTML}{D5DCE5}
\lstset{language=Python,basicstyle=\ttfamily\small,backgroundcolor=\color{codebg},frame=single,rulecolor=\color{codeframe},breaklines=true,showstringspaces=false,columns=fullflexible,keepspaces=true}
\IfFileExists{src/laptop.py}{}{\lstset{inputpath=../}}
\setlength{\parindent}{0pt}\setlength{\parskip}{5pt}

\begin{document}

\section{Scope and Data Flow}
This document contains only the dynamic-obstacle pipeline implemented in the current codebase:
\[
\text{LiDAR cluster}\rightarrow\text{track association}\rightarrow\text{obstacle EKF}\rightarrow\text{TCPA/DCPA}\rightarrow\text{virtual fields}\rightarrow\text{APF}.
\]
The vessel six-state EKF, thruster dynamics, COLREG decision details, route controller, GUI, and unrelated source files are excluded.

\section{Obstacle EKF}

\subsection{State and Constant-Velocity Model}
Each obstacle track uses
\[
\bm{x}_k=\begin{bmatrix}p_N&p_E&v_N&v_E\end{bmatrix}^{\mathsf T},
\]
where \(p_N,p_E\) are North--East position components and \(v_N,v_E\) are the corresponding velocities. The transition model is
\[
\bm{x}_{k+1}=\bm{F}(\Delta t)\bm{x}_k,
\qquad
\bm{F}=\begin{bmatrix}1&0&\Delta t&0\\0&1&0&\Delta t\\0&0&1&0\\0&0&0&1\end{bmatrix}.
\]
Position is advanced by velocity multiplied by the time step, while velocity is copied forward.

\subsection{Process Noise}
The constant-velocity assumption is imperfect because a real obstacle may accelerate or turn. The process covariance models this uncertainty:
\[
\bm{Q}=q\begin{bmatrix}
\frac{\Delta t^4}{4}&0&\frac{\Delta t^3}{2}&0\\
0&\frac{\Delta t^4}{4}&0&\frac{\Delta t^3}{2}\\
\frac{\Delta t^3}{2}&0&\Delta t^2&0\\
0&\frac{\Delta t^3}{2}&0&\Delta t^2
\end{bmatrix},\qquad q=\sigma_a^2.
\]
Larger \(\sigma_a\) makes the filter less confident in constant velocity and more responsive to measurements.

\lstinputlisting[firstline=1353,lastline=1387]{src/laptop.py}

\subsection{Prediction and Position Update}
The prediction is
\[
\hat{\bm{x}}_k=\bm{F}\bm{x}_{k-1}^{+},qquad
\hat{\bm{P}}_k=\bm{F}\bm{P}_{k-1}^{+}\bm{F}^{\mathsf T}+\bm{Q}.
\]
LiDAR measures position only:
\[
\bm{z}=\begin{bmatrix}p_N^{\mathrm{meas}}&p_E^{\mathrm{meas}}\end{bmatrix}^{\mathsf T},qquad
\bm{H}=\begin{bmatrix}1&0&0&0\\0&1&0&0\end{bmatrix}.
\]
The innovation, gain, corrected state, and Joseph covariance update are
\begin{align}
\bm{\nu}&=\bm{z}-\bm{H}\hat{\bm{x}},\\
\bm{S}&=\bm{H}\hat{\bm{P}}\bm{H}^{\mathsf T}+\bm{R},\\
\bm{K}&=\hat{\bm{P}}\bm{H}^{\mathsf T}\bm{S}^{-1},\\
\bm{x}^{+}&=\hat{\bm{x}}+\bm{K}\bm{\nu},\\
\bm{P}^{+}&=(\bm{I}-\bm{K}\bm{H})\hat{\bm{P}}(\bm{I}-\bm{K}\bm{H})^{\mathsf T}+\bm{K}\bm{R}\bm{K}^{\mathsf T}.
\end{align}
The covariance expresses uncertainty: the state is the best estimate, while the covariance determines how strongly the next measurement should correct it.

\lstinputlisting[firstline=1389,lastline=1459]{src/laptop.py}

\subsection{Track Stability and Future Position}
A track is not accepted as a reliable moving obstacle after a single detection. The implementation checks sample count, hit count, time span, speed, and missed detections. Stable tracks are extrapolated as
\[
\bm{p}(\tau)=\bm{p}_0+\bm{v}_0\tau+\frac12\bm{a}\tau^2.
\]
When acceleration is unavailable or too uncertain, \(\bm{a}=\bm{0}\) and the model becomes constant velocity.

\lstinputlisting[firstline=1461,lastline=1534]{src/laptop.py}
\lstinputlisting[firstline=1910,lastline=2118]{src/laptop.py}

\section{TCPA and DCPA}

Let own-ship position and velocity be \((\bm{p}_o,\bm{v}_o)\), and obstacle position and velocity be \((\bm{p}_t,\bm{v}_t)\). The implementation assumes that both vehicles keep their current velocities over the short prediction interval.

The relative position and velocity are
\[
\bm{r}=\bm{p}_t-\bm{p}_o,\qquad \bm{v}_r=\bm{v}_t-\bm{v}_o.
\]
The squared separation at time \(t\) is
\[
d^2(t)=\|\bm{r}+\bm{v}_r t\|^2
\]
and therefore
\[
d^2(t)=\bm{r}^{\mathsf T}\bm{r}+2t\bm{r}^{\mathsf T}\bm{v}_r+t^2\bm{v}_r^{\mathsf T}\bm{v}_r.
\]
Minimising this quadratic function gives the unconstrained closest-approach time
\[
t^*=-\frac{\bm{r}^{\mathsf T}\bm{v}_r}{\|\bm{v}_r\|^2}.
\]

\subsection{How the Sign Determines Closing or Separating}
The key quantity is \(\bm{r}^{\mathsf T}\bm{v}_r\):
\begin{itemize}
  \item If it is negative, the vehicles are closing and \(t^*>0\); the closest point is in the future.
  \item If it is positive, the vehicles are separating and \(t^*<0\); the mathematical closest point is in the past, so the current instant is used.
  \item If it is zero, relative motion is perpendicular to the current separation and the current instant is the closest point under the straight-line model.
  \item If \(\|\bm{v}_r\|^2\) is nearly zero, there is no reliable relative-motion direction; the implementation avoids division by a nearly zero value.
\end{itemize}

The implementation clips the non-negative time to the configured horizon:
\[
t_{\mathrm{CPA}}=\min\left(\max(t^*,0),T_{\max}\right).
\]
Thus, if the unrestricted closest point lies beyond the horizon, the reported distance is the distance at the horizon rather than an unrestricted long-term DCPA.

The CPA positions and distance are
\[
\bm{p}_{o,\mathrm{CPA}}=\bm{p}_o+\bm{v}_o t_{\mathrm{CPA}},
\quad
\bm{p}_{t,\mathrm{CPA}}=\bm{p}_t+\bm{v}_t t_{\mathrm{CPA}},
\]
\[
d_{\mathrm{CPA}}=\|\bm{p}_{t,\mathrm{CPA}}-\bm{p}_{o,\mathrm{CPA}}\|.
\]
TCPA answers ``when will the separation be smallest?''; DCPA answers ``how large will that separation be at the reported time?''.

\subsection{Numerical Examples}
For \(\bm{r}=[10,-5]^{\mathsf T}\) and \(\bm{v}_r=[-1,1]^{\mathsf T}\), the dot product is \(-15\) and the relative-speed square is \(2\). Therefore \(t^*=7.5\,\mathrm{s}\), so a \(10\,\mathrm{s}\) horizon returns \(t_{\mathrm{CPA}}=7.5\,\mathrm{s}\).

For \(\bm{r}=[10,0]^{\mathsf T}\) and \(\bm{v}_r=[1,0]^{\mathsf T}\), \(t^*=-10\,\mathrm{s}\). The code clips it to zero and reports the current separation as DCPA.

\lstinputlisting[firstline=60,lastline=85]{src/colreg_apf.py}
\lstinputlisting[firstline=2272,lastline=2278]{src/laptop.py}

\section{TCPA-Gated Virtual APF Fields}

Only a stable, sufficiently fast track satisfying
\[
0<t_{\mathrm{CPA}}\leq T_{\mathrm{horizon}}
\]
is converted into a virtual field. The final field centre is placed at
\[
\bm{p}_{\mathrm{virtual}}=\bm{p}_t+2t_{\mathrm{CPA}}\bm{v}_t.
\]
The factor two places the predicted field ahead of the closest-approach point and gives the planner time to react; it is an implementation rule, not a physical law.

\subsection{Own-Ship Equivalent Radius}
The obstacle is represented by an ellipse, while planning uses the own-ship reference point. The field therefore expands the obstacle semi-axes:
\[
a=\frac{pc1}{2}+r_{\mathrm{own}},\qquad b=\frac{pc2}{2}+r_{\mathrm{own}}.
\]
This approximates the own ship as a circle and inflates the obstacle footprint. Without this margin, the own-ship reference point could clear an obstacle while the physical hull still intersects it.

\subsection{Bridge Fields}
Intermediate ellipses are inserted between the current predicted position and the final \(2t_{\mathrm{CPA}}\) position. Their spacing is based on the combined support radius, producing a continuous predicted exclusion corridor instead of isolated future points.

\lstinputlisting[firstline=3745,lastline=3821]{src/laptop.py}

\section{APF Generation}

\subsection{Ellipse Geometry and Repulsion}
For obstacle centre \(\bm{c}\), length axis \(\hat{\bm{l}}\), width axis \(\hat{\bm{w}}\), and grid point \(\bm{p}\), define
\[
\ell=(\bm{p}-\bm{c})\cdot\hat{\bm{l}},\qquad
w=(\bm{p}-\bm{c})\cdot\hat{\bm{w}},
\]
\[
\rho=\sqrt{\left(\frac{\ell}{a}\right)^2+\left(\frac{w}{b}\right)^2}.
\]
The real-time field is active inside its outer scale \(s\). Between \(\rho=1\) and \(\rho=s\), the cubic fade is
\[
u=\operatorname{clip}\left(\frac{\rho-1}{s-1},0,1\right),qquad
f_{\mathrm{fade}}=1-(3u^2-2u^3),
\]
and
\[
\bm{F}_{\mathrm{rep}}=k_{\mathrm{rep}}f_{\mathrm{fade}}\hat{\bm{d}},
\]
where \(\hat{\bm{d}}\) points away from the obstacle.

\lstinputlisting[firstline=3618,lastline=3647]{src/laptop.py}
\lstinputlisting[firstline=200,lastline=208]{src/colreg_apf.py}
\lstinputlisting[firstline=3649,lastline=3711]{src/laptop.py}

\subsection{Attraction, Current Fields, and Virtual Fields}
The snapshot APF is the sum of goal attraction, current-obstacle fields, and risk-active virtual fields:
\[
U(\bm{p})=U_{\mathrm{goal}}(\bm{p})+
\sum_iU_{\mathrm{current},i}(\bm{p})+
\sum_jU_{\mathrm{virtual},j}(\bm{p}).
\]
The goal potential is
\[
U_{\mathrm{goal}}(\bm{p})=\frac12k_g\|\bm{p}-\bm{p}_{\mathrm{goal}}\|^2.
\]
The real-time controller uses vector repulsion; the snapshot renderer uses scalar grid potentials. Both use the same obstacle centres, dimensions, axes, and virtual-field records.

\lstinputlisting[firstline=770,lastline=841]{src/plot_apf_snapshots.py}
\lstinputlisting[firstline=897,lastline=950]{src/plot_apf_snapshots.py}

\section{Minimal Code Map and Verification}
\begin{longtable}{>{\raggedright\arraybackslash}p{0.30\textwidth}p{0.56\textwidth}}
\toprule
\textbf{Stage} & \textbf{Current implementation} \\
\midrule
\endfirsthead
\toprule
\textbf{Stage} & \textbf{Current implementation} \\
\midrule
\endhead
EKF prediction & \texttt{obstacle\_ekf\_predict} \\
Process noise & \texttt{obstacle\_ekf\_process\_noise} \\
Measurement update & \texttt{obstacle\_ekf\_update} \\
Stability and extrapolation & \texttt{obstacle\_track\_motion\_is\_stable}, \texttt{obstacle\_track\_state\_at} \\
Track association & \texttt{update\_apf\_obstacle\_tracks} \\
TCPA/DCPA & \texttt{straight\_line\_cpa}, \texttt{apf\_cpa\_metrics} \\
Virtual fields & \texttt{update\_apf\_virtual\_obstacles} \\
Real-time APF & \texttt{apf\_repulsion\_for\_obstacle} \\
Snapshot APF & \texttt{potential\_components} \\
\bottomrule
\end{longtable}

The focused regression checks are in \texttt{src/test\_colreg\_geometry.py}:
\begin{lstlisting}[language=bash]
cd src
python -m pytest test_colreg_geometry.py -q
\end{lstlisting}

\end{document}
